% =============================================================
% Block 2 -- Concepts and vocabulary (8 slides)
% Agent definition, chatbot vs agent, loop, prompt vs task,
% context/tokens, tools/permissions/sandbox, autonomy spectrum,
% reliability vocab.
% =============================================================

\sectionframe[accent]{Part 2 -- Concepts and vocabulary}

% -------- Slide 1: What is an agent? --------
\begin{frame}{What is an agent?}
  \begin{columns}[T,onlytextwidth]
    \column{0.46\textwidth}
      \bigpoint{LLM + tools\\+ goal + loop}
      \vspace{0.5cm}
      \begin{center}\color{muted}\small An agent is not a smarter chatbot.\\It is a workflow with partial autonomy.\end{center}
    \column{0.50\textwidth}
      \begin{tikzpicture}[every node/.style={align=center,font=\small}]
        \node[draw=deepblue,fill=softblue,rounded corners,very thick,
              minimum width=2.4cm,minimum height=0.9cm]
              (llm){\textbf{LLM}\\\scriptsize reasoning core};
        \node[draw=accent,fill=softorange,rounded corners,very thick,
              minimum width=2.4cm,minimum height=0.9cm,right=0.6cm of llm]
              (tools){\textbf{Tools}\\\scriptsize browser, files, shell};
        \node[draw=deepblue,fill=softgreen,rounded corners,very thick,
              minimum width=2.4cm,minimum height=0.9cm,below=0.6cm of llm]
              (goal){\textbf{Goal}\\\scriptsize the user's task};
        \node[draw=accent,fill=softpurple,rounded corners,very thick,
              minimum width=2.4cm,minimum height=0.9cm,below=0.6cm of tools]
              (loop){\textbf{Loop}\\\scriptsize plan/act/check};
        \draw[<->,thick,color=accent] (llm.east)--(tools.west);
        \draw[<->,thick,color=accent] (goal.east)--(loop.west);
        \draw[<->,thick,color=deepblue,dashed] (llm.south)--(goal.north);
        \draw[<->,thick,color=deepblue,dashed] (tools.south)--(loop.north);
      \end{tikzpicture}
  \end{columns}
\end{frame}

% -------- Slide 2: Chatbot vs agent --------
\begin{frame}{Chatbot vs agent}
  \vspace{0.2cm}
  \begin{center}
  \begin{tabularx}{0.92\linewidth}{>{\bfseries\color{deepblue}}lXX}
  \toprule
   & Chatbot & Agent\\
  \midrule
  Main action & Responds & Plans and acts\\
  Tools & Optional & Central\\
  Files & Usually attached & Often browsed and edited\\
  Failure mode & Wrong \emph{answer} & Wrong \emph{action}\\
  Human role & Reader & Supervisor\\
  Cost driver & Tokens read/written & Tokens \emph{plus} tool calls and retries\\
  \bottomrule
  \end{tabularx}
  \end{center}
  \vspace{0.4cm}
  \midpoint{Different failure mode $\Rightarrow$ different control needs.}
\end{frame}

% -------- Slide 3: Agent loop with a real trace --------
\begin{frame}{The agent loop, with a real trace}
  \begin{columns}[T,onlytextwidth]
    \column{0.36\textwidth}
      \begin{tikzpicture}[every node/.style={align=center,font=\scriptsize}]
        \foreach \i/\ang/\txt in {0/90/Observe,1/18/Plan,2/-54/Act,3/-126/Verify,4/162/Iterate}{
          \node[draw=deepblue,fill=softblue,rounded corners,very thick,
                minimum width=1.9cm,minimum height=0.65cm] at (\ang:2.1) (n\i) {\txt};}
        \draw[->,line width=1.2pt,color=accent] (n0) to[bend left=10] (n1);
        \draw[->,line width=1.2pt,color=accent] (n1) to[bend left=10] (n2);
        \draw[->,line width=1.2pt,color=accent] (n2) to[bend left=10] (n3);
        \draw[->,line width=1.2pt,color=accent] (n3) to[bend left=10] (n4);
        \draw[->,line width=1.2pt,color=accent] (n4) to[bend left=10] (n0);
        \node[font=\bfseries\small,color=deepblue] at (0,0) {Task};
      \end{tikzpicture}
    \column{0.60\textwidth}
      \scriptsize\ttfamily
      \begin{tcolorbox}[colback=softgray,colframe=deepblue!40,
                        arc=2pt,boxrule=0.4pt,top=2pt,bottom=2pt,left=4pt,right=4pt]
      \textbf{Task}: ``Re-run Table 2, fix the missing standard errors.''\\[-0.1cm]\hrule\vspace{0.1cm}
      \textbf{Observe}: read \texttt{do/02\_table2.do}, \texttt{tables/T2.tex}.\\
      \textbf{Plan}: \texttt{reghdfe} call lacks \texttt{cluster(id\_firm)}.\\
      \textbf{Act}: edit line 47, re-run \texttt{stata-mp -b do/02\_table2.do}.\\
      \textbf{Verify}: log shows clustered SE; \texttt{T2.tex} regenerated, columns 1--3 unchanged within 1e-4.\\
      \textbf{Iterate}: report diff to user; await approval.
      \end{tcolorbox}
  \end{columns}
  \smallsource{Stylised trace from a real Stata project; commands and file names anonymised.}
\end{frame}

% -------- Slide 4: From prompt to task --------
\begin{frame}{From prompt to task}
  \begin{columns}[T,onlytextwidth]
    \column{0.46\textwidth}
      \textbf{Prompt}\\[0.2cm]
      ``Explain how to run this regression.''
      \vspace{0.4cm}\\
      \scriptsize\color{muted}One question, one answer. No state. No side effect.
    \column{0.50\textwidth}
      \textbf{Task}\\[0.2cm]
      ``Inspect the scripts, re-run the model with clustered SE, regenerate the table, compare with the previous version, report what changed.''
      \vspace{0.4cm}\\
      \scriptsize\color{muted}A goal, a workflow, a verifiable artefact.
  \end{columns}
  \vspace{0.5cm}
  \bigpoint{Agentic AI shifts\\the unit of delegation.}
\end{frame}

% -------- Slide 5: Context, tokens, memory --------
\begin{frame}{Context, tokens, memory}
  \begin{columns}[T,onlytextwidth]
    \column{0.50\textwidth}
      \begin{tikzpicture}
        \begin{axis}[
            width=6.6cm,height=4.4cm,
            xbar,bar width=0.30cm,
            xmin=0,xmax=2200,
            xtick={0,200,500,1000,2000},
            xticklabels={0,200k,500k,1M,2M},
            xticklabel style={font=\scriptsize,color=muted},
            symbolic y coords={GPT-4 2023,Claude 3.5 2024,GPT-5 2025,Gemini 2.5 2026},
            ytick=data,
            yticklabel style={font=\scriptsize,color=dark},
            xlabel={Context window (tokens, approx.)},
            xlabel style={font=\scriptsize,color=muted},
            xmajorgrids,grid style={dotted,gray!40},
            axis lines*=left,
            tick style={color=muted},
            nodes near coords,
            nodes near coords style={font=\scriptsize\bfseries,color=deepblue,/pgf/number format/.cd,fixed,precision=0},
            point meta=explicit symbolic,
          ]
          \addplot[draw=deepblue,fill=deepblue!75]
            coordinates {(8,GPT-4 2023)[8k] (200,Claude 3.5 2024)[200k]
                         (400,GPT-5 2025)[400k] (2000,Gemini 2.5 2026)[2M]};
        \end{axis}
      \end{tikzpicture}
    \column{0.46\textwidth}
      \threebullets%
        {\textbf{Context} = what the model can see right now (prompt + files + tool outputs).}%
        {\textbf{Tokens} are the unit of measurement \emph{and} of cost.}%
        {\textbf{Memory} extends context across sessions (project notes, vault, vector store).}
      \vspace{0.2cm}
      \begin{center}\small\color{accent}\bfseries More context is not the same as better reasoning.\end{center}
  \end{columns}
  \smallsource{Vendor documentation 2023--2026 (OpenAI, Anthropic, Google). Numbers are nominal context-window sizes; effective context can be smaller.}
\end{frame}

% -------- Slide 6: Tools, permissions, sandbox --------
\begin{frame}{Tools, permissions, sandbox}
  \begin{center}
  \begin{tikzpicture}[every node/.style={align=center,font=\small}]
    \node[draw=deepblue,fill=softgreen,rounded corners,very thick,
          minimum width=3.2cm,minimum height=1.0cm,text width=3cm]
          (read){\textbf{read only}\\\scriptsize browse, summarise};
    \node[draw=deepblue,fill=softblue,rounded corners,very thick,
          minimum width=3.2cm,minimum height=1.0cm,text width=3cm,right=0.6cm of read]
          (edit){\textbf{edit workspace}\\\scriptsize change files};
    \node[draw=accent,fill=softorange,rounded corners,very thick,
          minimum width=3.2cm,minimum height=1.0cm,text width=3cm,right=0.6cm of edit]
          (run){\textbf{run commands}\\\scriptsize shell, scripts};
    \node[draw=accent,fill=softpurple,rounded corners,very thick,
          minimum width=3.2cm,minimum height=1.0cm,text width=3cm,right=0.6cm of run]
          (net){\textbf{network / API}\\\scriptsize external calls};
    \draw[->,thick,color=muted] (read)--(edit);
    \draw[->,thick,color=muted] (edit)--(run);
    \draw[->,thick,color=muted] (run)--(net);
  \end{tikzpicture}
  \end{center}
  \vspace{0.4cm}
  \begin{itemize}\setlength\itemsep{0.3em}
    \item \textbf{Tool} = something the model can call: browser, shell, files, database, an MCP server.
    \item \textbf{Permission} = the boundary between suggestion and action.
    \item \textbf{Sandbox} = a controlled environment where mistakes are reversible.
  \end{itemize}
  \midpoint{Permissions are the practical meaning of control.}
\end{frame}

% -------- Slide 7: Autonomy spectrum --------
\begin{frame}{The autonomy spectrum}
  \begin{center}
  \begin{tikzpicture}[every node/.style={align=center,font=\small}]
    \node[draw=deepblue,fill=softgreen,rounded corners,very thick,
          minimum width=2.7cm,minimum height=1.1cm,text width=2.5cm]
          (sug){\textbf{Suggest}\\\scriptsize human types it};
    \node[draw=deepblue,fill=softblue,rounded corners,very thick,
          minimum width=2.7cm,minimum height=1.1cm,text width=2.5cm,right=0.6cm of sug]
          (rev){\textbf{Review}\\\scriptsize human accepts diff};
    \node[draw=accent,fill=softorange,rounded corners,very thick,
          minimum width=2.7cm,minimum height=1.1cm,text width=2.5cm,right=0.6cm of rev]
          (exe){\textbf{Execute}\\\scriptsize agent acts in scope};
    \node[draw=accent,fill=softpurple,rounded corners,very thick,
          minimum width=2.7cm,minimum height=1.1cm,text width=2.5cm,right=0.6cm of exe]
          (orc){\textbf{Orchestrate}\\\scriptsize agent runs other agents};
    \draw[->,line width=1.4pt,color=accent] (sug)--(rev);
    \draw[->,line width=1.4pt,color=accent] (rev)--(exe);
    \draw[->,line width=1.4pt,color=accent] (exe)--(orc);
    \node[anchor=north,below=0.35cm of sug,color=muted,font=\scriptsize] {low risk};
    \node[anchor=north,below=0.35cm of orc,color=muted,font=\scriptsize] {high risk};
  \end{tikzpicture}
  \end{center}
  \vspace{0.5cm}
  \begin{center}\large\color{deepblue}\bfseries Pick the lowest level that gets the job done.\end{center}
  \vspace{0.2cm}
  \begin{center}\small\color{muted}For research, ``review'' and ``execute'' cover almost everything; ``orchestrate'' is for power users.\end{center}
\end{frame}

% -------- Slide 8: Reliability vocabulary --------
\begin{frame}{Reliability vocabulary you will keep hearing}
  \vspace{0.2cm}
  \begin{center}
  \begin{tabularx}{0.92\linewidth}{>{\bfseries\color{deepblue}}lX}
  \toprule
  Hallucination & A confident output that does not correspond to reality (a citation, a number, a function).\\
  Drift & The model's outputs change over time as providers update underlying weights or settings.\\
  Automation bias & The human tendency to trust automated outputs more than they deserve \citep{skitka_1999_automation_bias}.\\
  Reward hacking & The agent finds a shortcut that satisfies the metric but defeats the goal.\\
  Sycophancy & The model agrees with the user even when the user is wrong.\\
  \bottomrule
  \end{tabularx}
  \end{center}
  \vspace{0.4cm}
  \midpoint{These are research-process risks, not just AI quirks.}
\end{frame}
